\section{INTRODUCTION}\label{sec:intro}

Data pipeline failures are easy to describe and costly to fix: an upstream schema changes type, a
source stalls, a co-tenant starves compute, or ingress bursts past capacity. An empirical study of
real pipelines catalogs 41 distinct factors behind such quality and reliability
failures~\cite{foidl2024datapipeline}. The usual remedy is an on-call human, and that response time
dominates recovery time. Large language model (LLM) agents could cut that loop short, and systems
that apply them to cloud incidents are already running in production~\cite{chen2024rcacopilot,wang2023rcagent,luo2025opsagent}.
Most of those systems only recommend a diagnosis. Once an agent is allowed to \emph{act} on
production infrastructure, a governance question follows: what bounds what it can do?

Kirubakaran et~al.~\cite{kirubakaran2025governing} answer with four bounded agents, covering
monitoring, optimization, schema, and recovery, each producing only a structured action proposal
that a policy engine gates before anything executes. They report reductions of up to 45\% in
recovery time, about 25\% in cost, and over 70\% in manual interventions against static
orchestration. Four gaps separate that claim from the evidence for it.
\begin{enumerate}
  \item \emph{Evaluation path.} The original study names several model families as backends but gives too little implementation detail to pin down the exact model-in-the-loop condition behind its headline figures, even though work on the sim-to-real gap finds that simulated evaluation environments can inflate agent performance~\cite{liu2026simtoreal,zhu2026realusersim,gupta2026reliabilitybench}.
  \item \emph{Baselines.} The only comparator is static orchestration: fixed allocations, predefined retries, human-initiated remediation, and autoscaling disabled, despite mature non-LLM automation already running at production scale~\cite{rzadca2020autopilot}.
  \item \emph{Adversarial testing.} The safety case rests on the policy gate, yet a stated bound is not the same as a tested one. Agent-security benchmarks report high attack success against weak
    defenses~\cite{debenedetti2024agentdojo,zhang2024asb}, and the policy language itself has been
    criticized as too weak for agentic governance~\cite{joshi2026deontic}.
  \item \emph{Statistics and generality.} No confidence intervals, effect sizes, or run counts are reported. The GPT, Claude, and Gemini model families are described as interchangeable backends whose differences touched only verbosity and latency, yet no cross-model results back that claim, and agent behavior has elsewhere been shown to vary across models and providers~\cite{sanna2025crossllm,zhang2026towardagentic}.
\end{enumerate}
We build the system and put five questions to it.
\begin{enumerate}
  \item[\textbf{RQ1:}] With a real model in the loop, does the full agent system outperform static
    orchestration on recovery time, cost, and manual interventions?
  \item[\textbf{RQ2:}] Does it outperform cheap, non-LLM automation, or only a slow human?
  \item[\textbf{RQ3:}] Does it act correctly, not merely quickly?
  \item[\textbf{RQ4:}] Do the answers to RQ1--RQ3 change once the model is swapped for a
    deterministic mock?
  \item[\textbf{RQ5:}] Does the policy gate contain adversarial proposals, judged against an oracle
    independent of the policy code?
\end{enumerate}

\paragraph{Central finding}
RQ4 proves decisive. Running the identical \texttt{full} configuration first under a deterministic
mock and then under a real, remote model produces results that disagree by roughly three orders of
magnitude on recovery time, and on whether the executed mitigation was even correct
(Section~\ref{sec:results}). The mock run is not wrong in any ordinary sense; it faithfully measures
a faster, more forgiving experiment. The trouble is that a governance claim drawn from that
experiment does not carry over to the one an operator actually faces. Every other number in this
paper is conditional on that divergence: an evaluation built on the mock's near-zero latency, or a
cost model that assumes away the provisioning gap it rests on, would report gains a live deployment
simply does not deliver.

\paragraph{Contributions}
A contribution counts here only if this campaign both implements and measures it. Four qualify.
\begin{enumerate}
  \item \textbf{A live-model evaluation} spanning eight configurations, namely three non-agent
    comparators, the full system, and four single-agent ablations, across four injected fault
    scenarios. The campaign totals \NRunsTotal{} runs, with per-run provenance and billing-accurate
    LLM accounting (Section~\ref{sec:method}).
  \item \textbf{Explicit non-LLM baselines and a decision-quality metric.} A rule-based responder and
    threshold autoscaling, each modelled with a fixed remediation latency and run through the same
    campaign, show whether the agent system beats cheap automation and not just static orchestration;
    a correct-mitigation metric shows whether it acts correctly, not merely quickly
    (Section~\ref{sec:results}).
  \item \textbf{An adversarial evaluation against an independent oracle.} A generated corpus is
    graded by a specification written separately from the policy code under test, so the resulting
    containment rate cannot be circular (Sections~\ref{sec:results} and~\ref{sec:lessons}).
  \item \textbf{A disclosed, testable cost model} that turns the base paper's loosely specified cost
    metric into a reproducible one, together with closed-form sensitivity analyses for the two
    assumptions the headline numbers depend on: the simulated human's latency and the cost model's
    provisioning gap (Eq.~\ref{eq:cost}).
\end{enumerate}
Two further contributions are implemented but this campaign never exercised them, so we present them
as architecture rather than results. A cross-model harness could test the base paper's model-agnostic
claim, but its live sweep has not yet run against real providers. A production trust core, with
graduated autonomy, a durable kill switch, and a per-target blast-radius cap, exists too, but every
run in the campaign used one fixed mode, leaving no mode-transition or kill-switch-latency
measurement to report. Section~\ref{sec:system} describes both, and Section~\ref{sec:discussion}
returns to them as the concrete gaps future work would close. Finally, a set of
\textbf{production-hardening lessons} (Section~\ref{sec:lessons}) records defects this system
actually exhibited.

\paragraph{What is not novel}
The architecture itself, its three planes, four bounded agents, policy gate, and
observe--reason--propose loop, belongs to Kirubakaran et~al.; we implemented it rather than invented
it. We claim no new agent design, policy language, or statistical method. What is ours is the
evidence, the measurement discipline that produced it, and the defects uncovered by holding the
system to its own stated guarantees rather than to a test suite built around the same assumptions the
system makes.

\paragraph{Organization}
Section~\ref{sec:related} surveys related work, Section~\ref{sec:problem} formalizes the problem and
the metrics, and Section~\ref{sec:system} describes the system as built. Section~\ref{sec:method}
lays out the experimental design, and Section~\ref{sec:results} presents the results around the
mock-versus-live divergence. Section~\ref{sec:lessons} gives the hardening lessons,
Section~\ref{sec:discussion} covers limitations and threats to validity, and
Section~\ref{sec:conclusion} concludes.
